\subsection{Contenido original e incidentes de integridad}\label{sec:impl-contenido-documental}

La consulta de contenido identifica expediente, UUID documental y versión positiva. Recupera los bytes originales tanto si están pendientes de sello como si ya contienen evidencia capturada; no sustituye la verificación de firma, certificado o sello de tiempo. Una versión nueva no redirige la petición de una anterior. Owner accede a todos los expedientes y Litigante o Paralegal requieren pertenencia vigente. Cliente continúa denegado; el cierre administrativo conserva esta lectura autorizada.

El servicio carga una captura cifrada autorizada, valida el sobre, descifra mediante AES-GCM con el AAD de identidad y versión y compara SHA-256. Después reautentica al principal completo. La transacción final comprueba cuenta, rol, pertenencia y coincidencia del contenido inmutable, y confirma el evento \texttt{document.content\_authorized}. Sólo entonces el transporte recibe los bytes. El evento acredita autorización confirmada, no recepción por el navegador ni escritura en disco. Redis y PostgreSQL no forman una transacción distribuida; cerrar sesión después de liberar los bytes no puede recuperarlos.

La ruta admite hasta 16~MiB de contenido. PostgreSQL comprueba el tamaño cifrado antes de materializarlo y repite la cota en la selección; el procesador aplica su propio límite. El cupo de descargas acompaña al buffer ceroizable hasta liberar la última referencia de bytes, incluso si el manejador HTTP ya terminó; un receptor lento no libera prematuramente esa capacidad. Un histórico mayor devuelve error de tamaño sin eliminarse ni generar un incidente de integridad. Consultar tampoco vuelve a admitir el formato de una carga histórica. La respuesta satisfactoria usa adjunto binario, nombre saneado, longitud exacta, prohibición de caché y cabeceras de expediente, identidad, versión y digest. No se entregan rangos parciales: la validación comprende el archivo completo.

Un fallo de validación produce una transacción independiente que conserva el incidente y \texttt{document.content\_rejected}; no crea un acceso satisfactorio. Registra identidades exactas, solicitante histórico, clase del fallo, instantes y huellas, sin texto claro, vault, tokens ni claves. El actor técnico identifica la observación y no acusa al solicitante. El UUID de observación permite repetir exactamente el registro sin duplicarlo; reutilizarlo con otros datos produce conflicto. Intentos distintos mantienen incidentes distintos.

Las clases separan sobre inválido, autenticación criptográfica fallida, digest discordante y cambio de captura. No determinan la causa: una KEK incorrecta también puede causar fallo de autenticación. Un error de base de datos o auditoría sigue siendo técnico; si no se confirma el incidente, no se afirma que Owner haya sido notificado. El rechazo conserva cero bytes de archivo en la respuesta.

La tabla \texttt{document\_integrity\_incidents} es inmutable y forma parte del catálogo, inventario y respaldo. Sólo Owner activo consulta su buzón, con autorización y auditoría actuales; el orden paginado es UUID ascendente, no cronológico. Qadra consulta existencia al entrar, permite actualizar y abre la versión exacta tras revalidar acceso. La apertura no descarga automáticamente contenido rechazado. El aviso no acredita correo, reparación ni atención humana. La descarga descarta resultados tardíos al cambiar expediente, versión o sesión. El contrato se conserva en \rutarepo{docs/document-content-api.md} y la decisión en el ADR~0042 de \rutarepo{docs/adr/}.
